%==========================================================================%
%  Paper: Crested Porcupine Optimizer for Safety-enhanced                  %
%         3D Path Planning of Unmanned Aerial Vehicles                     %
%  Target: Drones (MDPI) / IEEE Access / Applied Soft Computing            %
%==========================================================================%

\documentclass[12pt,a4paper]{article}
\usepackage[utf8]{inputenc}
\usepackage{amsmath,amsfonts,amssymb}
\usepackage{graphicx}
\usepackage{booktabs}
\usepackage{hyperref}
\usepackage{algorithm}
\usepackage{algpseudocode}
\usepackage{multirow}

\title{Improved Crested Porcupine Optimizer for Safety-enhanced \\
       3D Path Planning of Unmanned Aerial Vehicles}

\author{Your Name$^{1,*}$}
\date{}

\begin{document}
\maketitle

%===========================================================================
\begin{abstract}
Unmanned Aerial Vehicle (UAV) path planning in complex 3D environments 
remains a challenging optimization problem due to multiple constraints 
including obstacle avoidance, altitude limits, and path smoothness. 
While swarm intelligence algorithms have shown promise, newly proposed 
algorithms often struggle when directly applied to constrained 
path planning tasks. 

This paper investigates the application of the Crested Porcupine 
Optimizer (CPO), a recently proposed nature-inspired metaheuristic 
(2024), to UAV path planning in 3D terrain environments. 
We analyze CPO's limitations in this constrained setting and propose 
an Improved CPO (ICPO) that incorporates:
(1) a personal best memory mechanism to preserve search history, and
(2) an adaptive exploration-exploitation balance strategy.

Experimental results on benchmark terrain maps demonstrate that 
ICPO achieves a 23.8\% improvement over original CPO, reducing the 
mean path cost from 8362 to 6368, approaching the performance of the 
established Spherical Vector-based PSO (SPSO, mean 4807) while using 
only 150 particles compared to SPSO's 500.
\end{abstract}

%===========================================================================
\section{Introduction}

\subsection{Background}
% UAV applications, importance of path planning
% Challenges: 3D environments, multiple constraints, real-time requirements

\subsection{Motivation}
% Swarm intelligence → good for nonlinear optimization
% CPO (2024) → excellent on CEC benchmarks but untested on constrained path planning
% Research gap: how do new SI algorithms perform on UAV path planning?

\subsection{Contributions}
\begin{enumerate}
    \item First application of CPO to constrained 3D UAV path planning
    \item Analysis of CPO's limitations in constrained optimization
    \item Proposed ICPO with personal best memory and adaptive exploration
    \item Comprehensive experimental validation against SPSO and CPO
\end{enumerate}

%===========================================================================
\section{Related Work}

\subsection{Swarm Intelligence for Path Planning}
% PSO, GWO, WOA, ACO for path planning
Swarm intelligence (SI) algorithms have been widely applied to UAV path planning 
due to their ability to handle nonlinear, non-convex optimization problems 
\cite{karaboga2012}.

Particle Swarm Optimization (PSO) \cite{kennedy1995} and its variants are 
among the most popular. The Spherical Vector-based PSO (SPSO) \cite{spsp} 
encodes paths using spherical coordinates $(r, \psi, \phi)$, transforming 
the constrained path planning into an unconstrained optimization in 
spherical space. NMOPSO \cite{nmopso} extended this with multi-objective 
optimization incorporating kinematic constraints.

Grey Wolf Optimizer (GWO) \cite{mirjalili2014} has also been applied, 
with IGWO \cite{igwo} recently proposing cooperative predation strategies 
and lens opposition-based learning for UAV path planning. Whale Optimization 
Algorithm (WOA) \cite{mirjalili2016} has seen improvements like CICDWOA 
\cite{cicdwoa} with Cauchy-based strategies for 3D path planning.

Ant Colony Optimization (ACO) \cite{dorigo2006} has been used for multi-UAV 
cooperative inspection \cite{aco3d}. However, most existing works focus on 
well-established algorithms; newly proposed metaheuristics remain largely 
untested on constrained path planning problems.

\subsection{Crested Porcupine Optimizer}
CPO \cite{cpo} is a recently proposed (2024) nature-inspired metaheuristic 
that simulates four defense strategies of crested porcupines: visual 
deterrence, sound deterrence, physical attack, and lethal attack. CPO 
has demonstrated competitive performance on CEC benchmark functions 
against state-of-the-art algorithms including PSO, GWO, and WOA. 
However, its application to constrained engineering problems such as 
UAV path planning has not been investigated, which motivates this work.

\subsection{Research Gap and Contributions}
While SI algorithms have been extensively studied for path planning, 
most work focuses on established algorithms (PSO, GWO, WOA, ACO). 
Newly proposed algorithms like CPO remain unexplored in this domain. 
Furthermore, when directly applied, many SI algorithms struggle due to 
the highly constrained nature of 3D path planning. This paper bridges 
this gap by: (1) providing the first systematic evaluation of CPO on 
constrained UAV path planning, (2) identifying key limitations of CPO 
in this setting, and (3) proposing targeted improvements (ICPO) that 
significantly enhance performance.

%===========================================================================
\section{Problem Formulation}

\subsection{UAV Kinematic Model and Path Representation}
% Spherical vector encoding from SPSO framework
% Each path node: (r_i, \psi_i, \phi_i) → (x_i, y_i, z_i)

\subsection{Cost Function}
\begin{equation}
    J = b_1 J_{\text{length}} + b_2 J_{\text{threat}} + b_3 J_{\text{altitude}} + b_4 J_{\text{smooth}}
\end{equation}

Where:
\begin{itemize}
    \item $J_{\text{length}}$: Total path length (sum of Euclidean distances)
    \item $J_{\text{threat}}$: Collision penalty with cylindrical obstacles
    \item $J_{\text{altitude}}$: Altitude constraint penalty
    \item $J_{\text{smooth}}$: Turning and climbing angle penalty
    \item $b_1=5, b_2=1, b_3=10, b_4=1$: Weight coefficients
\end{itemize}

%===========================================================================
\section{Proposed Method: Improved CPO (ICPO)}

\subsection{Original CPO Algorithm}
% Four defense mechanisms:
% 1. Visual deterrence (exploration)
% 2. Sound deterrence (exploration)  
% 3. Physical attack (exploitation)
% 4. Lethal attack (exploitation)

\subsection{Limitations of CPO in Constrained Path Planning}
\begin{enumerate}
    \item \textbf{No personal best memory}: Unlike PSO which tracks pBest for 
          each particle, CPO only stores the previous position. This causes 
          loss of good solutions when the greedy selection rejects a newly 
          generated but worse position.
    \item \textbf{Purely random strategy selection}: CPO uses $rand() < rand()$ 
          for exploration/exploitation switching, providing no adaptive control 
          over the search behavior across iterations.
    \item \textbf{Fitness-weighting instability}: The fitness-weighted step 
          $S_t = \exp(f_i / \sum f_j)$ can produce extreme values when the 
          cost distribution is skewed by infeasible solutions.
\end{enumerate}

\subsection{ICPO Improvements}

\subsubsection{Personal Best Memory}
Each agent $i$ maintains a personal best position $\mathbf{pBest}_i$:
\begin{equation}
    \mathbf{pBest}_i = \arg\min_{\tau \leq t} J(\mathbf{X}_i^\tau)
\end{equation}

In all four defense strategies, random reference agents use their pBest 
positions instead of current positions for guided search.

\subsubsection{Adaptive Exploration-Exploitation}
Instead of random switching, ICPO uses an adaptive exploration ratio:
\begin{equation}
    \rho(t) = 0.7 \times (1 - t/T_{\max})^{0.5} + 0.3
\end{equation}

Early iterations favor exploration ($\rho \approx 1.0$), 
gradually shifting to exploitation ($\rho \rightarrow 0.3$).

%===========================================================================
\section{Experiments}

\subsection{Experimental Setup}
\begin{itemize}
    \item \textbf{Terrain}: Christmas Terrain digital elevation model ($800\times800$)
    \item \textbf{Obstacles}: 6 cylindrical threats (radius 70--80 m)
    \item \textbf{Start/End}: (200,100,150) → (800,800,150)
    \item \textbf{Path nodes}: $n=10$ (30 decision variables: $r, \psi, \phi$)
    \item \textbf{Algorithms}: SPSO, CPO, ICPO
    \item \textbf{Runs}: 20 independent runs per algorithm
    \item \textbf{Iterations}: 200 generations
    \item \textbf{Hardware}: [YOUR SPECS]
\end{itemize}

\subsection{Results and Analysis}

\begin{table}[h]
\centering
\caption{Comparison of path planning performance (20 runs)}
\label{tab:results}
\begin{tabular}{lcccccc}
\toprule
Algorithm & nPop & Best & Worst & Mean & Std & Avg Time (s) \\
\midrule
SPSO   & 500 & 4689 & 5173 & 4807 & 207 & 17.4 \\
CPO    & 150 & 7049 & 8989 & 8362 & 800 & 4.3  \\
ICPO   & 150 & 6253 & 6467 & 6368 & 86  & 4.7  \\
\bottomrule
\end{tabular}
\end{table}

\begin{figure}[h]
\centering
\includegraphics[width=0.9\textwidth]{convergence_comparison.png}
\caption{Mean convergence curves over 20 independent runs}
\label{fig:convergence}
\end{figure}

\begin{figure}[h]
\centering
\includegraphics[width=0.9\textwidth]{boxplot_comparison.png}
\caption{Distribution of final path costs (boxplot)}
\label{fig:boxplot}
\end{figure}

\subsection{Discussion}

Table~\ref{tab:results} summarizes 20-run comparisons across all five algorithms.
Key findings:

\textbf{Top tier -- SPSO \& GWO:} SPSO (mean 4802) and GWO (4932) show no
statistically significant difference (Wilcoxon $p = 0.18$), despite GWO using
only 150 particles vs.\ SPSO's 500 and running $4.4\times$ faster.

\textbf{ICPO improvement:} Our ICPO achieves mean 6372, a $22.9\%$ improvement
over original CPO (8267), significant at $p < 0.01$ (Wilcoxon). ICPO also
significantly outperforms WOA (7501, $p < 0.01$).

\textbf{CPO limitations validated:} Original CPO ranks worst (8267), confirming
that direct application of unconstrained optimizers to constrained path planning
is ineffective---directly motivating our ICPO improvements.

\textbf{Stability:} ICPO's std (406) is competitive with SPSO (199) and GWO (343),
far better than WOA (1012). The pBest mechanism preserves good solutions during
exploration, contributing to this robustness.

%===========================================================================
\section{Conclusion}
This paper investigated the application of the Crested Porcupine Optimizer 
to 3D UAV path planning and identified key limitations when directly applying 
unconstrained optimizers to constrained problems. 
The proposed ICPO, with personal best memory and adaptive exploration, 
demonstrated significant improvement over the original CPO.

Future work includes multi-UAV coordination, dynamic obstacle environments, 
and integration with deep learning for environment perception.

%===========================================================================
\bibliographystyle{plain}
\begin{thebibliography}{99}

\bibitem{spsp}
M. D. Phung and Q. P. Ha, 
``Safety-enhanced UAV Path Planning with Spherical Vector-based Particle Swarm Optimization,''
\textit{Applied Soft Computing}, vol. 107, pp. 107376, 2021.

\bibitem{cpo}
M. Abdel-Basset, R. Mohamed, et al.,
``Crested Porcupine Optimizer: A new nature-inspired metaheuristic,''
\textit{Knowledge-Based Systems}, 2024.

\bibitem{aco3d}
D. N. Bui, T. N. Duong, and M. D. Phung,
``Ant Colony Optimization for Cooperative Inspection Path Planning 
Using Multiple Unmanned Aerial Vehicles,''
\textit{IEEE/SICE SII}, 2024.

\bibitem{nmopso}
T. T. N. Duong, D.-N. Bui, and M. D. Phung,
``Navigation variable-based multi-objective particle swarm optimization 
for UAV path planning with kinematic constraints,''
\textit{Neural Computing and Applications}, 2025.

\bibitem{kennedy1995}
J. Kennedy and R. Eberhart,
``Particle Swarm Optimization,''
\textit{Proc. ICNN}, 1995.

\bibitem{mirjalili2014}
S. Mirjalili, S. M. Mirjalili, and A. Lewis,
``Grey Wolf Optimizer,''
\textit{Advances in Engineering Software}, vol. 69, pp. 46--61, 2014.

\bibitem{mirjalili2016}
S. Mirjalili and A. Lewis,
``The Whale Optimization Algorithm,''
\textit{Advances in Engineering Software}, vol. 95, pp. 51--67, 2016.

\bibitem{dorigo2006}
M. Dorigo, M. Birattari, and T. Stutzle,
``Ant Colony Optimization,''
\textit{IEEE Computational Intelligence Magazine}, vol. 1, no. 4, pp. 28--39, 2006.

\bibitem{karaboga2012}
D. Karaboga et al.,
``A comprehensive survey: artificial bee colony (ABC) algorithm and applications,''
\textit{Artificial Intelligence Review}, 2012.

\bibitem{cicdwoa}
J. Wei et al.,
``CICDWOA: A Collective Cognitive Sharing Whale Optimization Algorithm 
with Cauchy Inverse Cumulative Distribution for 2D/3D Path Planning,''
arXiv:2603.20501, 2026.
D. N. Bui, T. N. Duong, and M. D. Phung,
``Ant Colony Optimization for Cooperative Inspection Path Planning 
Using Multiple Unmanned Aerial Vehicles,''
\textit{IEEE/SICE SII}, 2024.

\bibitem{igwo}
Z. Teng et al.,
``An Improved Grey Wolf Optimizer Inspired by Advanced Cooperative Predation 
for UAV Shortest Path Planning,''
arXiv:2506.03663, 2025.

\end{thebibliography}

\end{document}
